\section{Off-policy 校正用权重交换分布偏差与方差}
\label{sec:16-Control-and-Reinforcement-Learning/02-Reinforcement-Learning/08}

线上跑着的推荐策略 \(\mu\) 带一点随机探索，半年攒下几千万条会话日志。产品经理拿来一个新策略 \(\pi\)，问你上线前能不能先估出它的日均收益。全量 A/B 要两周并且有收入风险，你手里只有 \(\mu\) 的数据。直接把日志里的平均回报报上去，回答的是另一个问题；于是你给每条会话乘上概率比再平均。前几次估计看着合理，某天一条会话的权重算出来是 \(1.2\times10^{4}\)，整个估计基本由它一条决定，去掉它结论反号。

权重把``这批数据来自谁''这件事写成了一个可计算的量。它同时把分布的偏差搬进了方差。

\subsection{比率里环境动力学自己消掉}

一段从 \(t\) 到 \(T\) 的轨迹，在目标策略下的概率是

\[
p_\pi(\tau)=\prod_{k=t}^{T-1}\pi(A_k\given S_k)\,P(S_{k+1}\given S_k,A_k),
\]

换成行为策略只需把 \(\pi\) 换成 \(\mu\)。两者的环境是同一个，转移概率逐项对消：

\[
W_{t:T-1}=\frac{p_\pi(\tau)}{p_\mu(\tau)}=\prod_{k=t}^{T-1}\rho_k,
\qquad
\rho_k=\frac{\pi(A_k\given S_k)}{\mu(A_k\given S_k)}.
\]

于是在可积条件下

\[
V^\pi(s)=\E_\mu\big[W_{t:T-1}G_t\given S_t=s\big].
\]

这是 \hyperref[sec:13-Machine-Learning/08-Sampling-and-Inference/01]{``Monte Carlo 用随机样本逼近期望''} 里那套重要性采样，换了一身衣服：提议分布是行为策略诱导的轨迹分布，目标分布是目标策略诱导的。好消息是转移概率不必知道；代价是必须知道当时那个动作被选中的概率。日志只记了动作没记 propensity，事后拿一个模型去反推 \(\mu(a\given s)\)，模型误差会原封不动进入权重。

支持条件也照搬过来：凡是 \(\pi(a\given s)>0\) 的地方都要有 \(\mu(a\given s)>0\)。行为策略从没执行过的动作没有反事实结果，任何权重都变不出零样本里的信息。

\subsection{乘积权重按轨迹长度指数集中}

单步比率温和，乘积照样能失控，原因在于乘法在对数尺度上是求和。看一个极端简化的例子：每步 \(\rho_k\) 以各一半的概率取 \(3/2\) 或 \(1/2\)，各步独立。均值 \(\E[\rho]=1\)，看上去无害。但

\[
\E[\log\rho]=\tfrac12\big(\log\tfrac32+\log\tfrac12\big)\approx-0.144,
\]

五十步之后 \(\log W\) 的均值是 \(-7.2\)，典型权重落在 \(e^{-7.2}\approx7.5\times10^{-4}\)。同时 \(\E[W]=1\) 一点没变——均值靠的是极少数取到大量 \(3/2\) 的轨迹。二阶矩更直白：\(\E[\rho^2]=1.25\)，于是 \(\E[W^2]=1.25^{50}\approx7\times10^{4}\)，标准差约 \(265\) 而均值是 \(1\)。

普通估计量

\[
\widehat V_{\mathrm{IS}}=\frac1n\sum_{i=1}^{n}W_iG_i
\]

在支持条件与矩条件下无偏，可无偏说的是无限次重复的平均。上面那个例子里，要让有效样本量达到十几，需要几十万条轨迹；你手上那一批数据里，绝大多数轨迹的权重接近零，估计的命运握在两三条轨迹手里。

\subsection{自归一化换来稳定，代价是有限样本偏差}

一个便宜的改进是把权重归一化：

\[
\widehat V_{\mathrm{SN}}=\frac{\sum_iW_iG_i}{\sum_iW_i}.
\]

它对权重的整体尺度不敏感，方差通常小很多，在大数定律成立时依然一致，但有限样本下有偏。选它还是选 \(\widehat V_{\mathrm{IS}}\)，是一次明码标价的偏差---方差交换，别把归一化之后的估计继续称作严格无偏。

集中度的常用读数是有效样本量

\[
N_{\mathrm{eff}}=\frac{\big(\sum_iW_i\big)^2}{\sum_iW_i^{2}},
\]

它回答的问题只有一个：估计是不是被少数几条样本包办了。它管不了轨迹之间是否独立、propensity 模型是否正确，更管不了覆盖是否充分——若全部样本都漏掉了目标策略最爱去的区域，权重可能个个平淡，\(N_{\mathrm{eff}}\) 看上去很健康，而估计整体偏在一边。\hyperref[sec:13-Machine-Learning/08-Sampling-and-Inference/05]{``粒子滤波用带权样本追踪非线性状态''} 里的权重塌缩是同一个现象在时间维度上的版本。

\subsection{只给该负责的奖励配权重}

把整条轨迹的比率乘到整个回报上，等于让第 \(40\) 步的动作随机性去影响第 \(3\) 步已经发生的奖励。因果顺序允许把它删掉：

\[
V^\pi(s)=\E_\mu\left[\sum_{k=t}^{T-1}\gamma^{k-t}\left(\prod_{j=t}^{k}\rho_j\right)R_{k+1}\;\middle|\;S_t=s\right].
\]

奖励 \(R_{k+1}\) 只乘上生成它之前那些动作的比率，不再替后面的动作背方差。这项化简与下一节里用 reward-to-go 删掉过去奖励属于同一种手法：把与被估计量无关的随机性从估计量里请出去，期望不变，方差下降。

对一步 TD，校正常常缩成单个因子：

\[
w_{t+1}=w_t+\alpha_t\,\rho_t\,\delta_t\,\nabla_w\widehat V(S_t;w_t).
\]

要看清它做了什么、没做什么。\(\rho_t\) 把给定状态下的动作分布从 \(\mu\) 改成了 \(\pi\)；样本状态仍然按 \(d^\mu\) 出现，投影仍然是按 \(d^\mu\) 加权的投影。换句话说，前面那节里让复合算子失去压缩性的正是这个状态加权，而单步权重没碰它。重要性采样不拆 Deadly Triad，它处理的是三角里``动作''这一半。

表格 Q 学习不显式乘 \(\pi/\mu\)，靠 \(\max\) 算子直接逼近最优 Bellman 方程，前提是关键状态动作对都被访问够。一旦进入函数逼近或固定日志，缺失的覆盖会以外推误差的形式回来，只是不再挂着``权重''这个名字。

\subsection{截断是在偏差与方差之间选一个点}

工程上普遍的做法是给权重封顶，单步用 \(\bar\rho_t=\min(c,\rho_t)\)，或者给累积乘积设上限。封顶之后没有哪条轨迹能造成任意大的更新，而期望也随之被改动了。

\begin{center}
\begin{tikzpicture}[>=stealth,scale=0.85]
\draw[->] (0,0) -- (5.4,0) node[below,font=\small] {截断阈值 \(c\)};
\draw[->] (0,0) -- (0,3.4) node[left,font=\small] {误差};
\draw[dashed] plot[domain=0:3.4,samples=40] ({\x*1.4},{2.4*exp(-0.9*\x)});
\draw[dotted,thick] plot[domain=0:3.4,samples=40] ({\x*1.4},{0.225*exp(0.75*\x)});
\draw[thick] plot[domain=0:3.4,samples=40]
  ({\x*1.4},{2.4*exp(-0.9*\x)+0.225*exp(0.75*\x)});
\filldraw (2.31,1.31) circle (2pt);
\draw[dashed] (2.31,0) -- (2.31,1.31);
\node[below,font=\small] at (2.31,0) {选点};
\node[right,font=\small] at (2.1,2.9) {总误差};
\node[right,font=\small] at (0.75,1.1) {偏差};
\node[right,font=\small] at (3.2,0.72) {方差};
\begin{scope}[shift={(7.0,0)}]
\draw[->] (0,0) -- (5.0,0) node[below,font=\small] {轨迹};
\draw[->] (0,0) -- (0,3.4) node[left,font=\small] {\(W_i\)};
\foreach \x/\h in {0.3/0.06,0.65/0.04,1.0/0.10,1.35/0.03,1.7/2.95,2.05/0.05,
  2.4/0.07,2.75/0.03,3.1/1.10,3.45/0.04,3.8/0.06,4.15/0.02,4.5/0.05}
  \draw[very thick] (\x,0) -- (\x,\h);
\draw[dashed] (0,0.9) -- (4.8,0.9) node[right,font=\small] {\(c\)};
\node[right,font=\small,align=left] at (0.2,2.4) {\(N_{\mathrm{eff}}\approx2\)};
\end{scope}
\end{tikzpicture}
\end{center}

\emph{右边是没截断时权重的常态：几乎全部接近零，两条包办估计；左边是把顶封在 \(c\) 之后，偏差与方差沿相反方向走，你要挑的是那个总误差最低的点。}

阈值 \(c\) 小则稳而偏，大则准而抖，图上那个最低点的位置取决于轨迹长度、两个策略的差距和样本量，没有通用取值。要紧的是别把截断当成免费的补丁：报告里应当写清比率的分位数、被截断的样本比例，以及结论对阈值的敏感性。

V-trace 和 Retrace 比裸截断多做了一步，它们把截断放进一个新的多步算子里，让截断后的迭代仍有明确的不动点。V-trace 用两个不同的阈值：一个乘在时序差分项上，控制信息沿轨迹传播的强度和方差；另一个决定不动点对应哪个策略——阈值收紧时，那个不动点从 \(\pi\) 的价值滑向 \(\mu\) 与 \(\pi\) 之间某个策略的价值。知道自己算出来的是哪个策略的价值，比算出一个方差小但不知道对应谁的数强。

策略优化里的比率裁剪外形相似，目的不同：那里限制的是新策略相对于生成这批数据的策略的偏移幅度，属于信赖域式的约束。见到 \(\pi/\mu\) 就假定三者同源，会把策略约束、离策略评估和方差缩减混成一锅。

\subsection{在线与离线是另一条轴}

\begin{longtable}{p{0.16\textwidth}p{0.36\textwidth}p{0.36\textwidth}}
\toprule
 & On-policy & Off-policy \\
\midrule
在线 &
用当前 \(\pi\) 交互并立刻更新；数据最新鲜，但每更新一次，旧样本就成了别人的数据 &
用探索策略、旧策略或多个智能体持续收集，同时学习目标 \(\pi\)；回放缓冲属于这一格 \\
离线 &
固定日志恰好来自待评估的那个 \(\pi\)；能评估，不能探索 &
固定日志来自 \(\mu\)，要评估或改进另一个 \(\pi\)；覆盖与 propensity 成为硬约束 \\
\bottomrule
\end{longtable}

带回放的 DQN 在``交互还在进行''的同时是离策略的，这两件事被同一个词组绑在一起讨论时最容易混淆。纯离线的场景更紧：支持缺口补不上，通常还要把新策略约束在数据附近，或对分布外动作施加悲观估计。

\subsection{有些时候该承认证据不够}

最大权重随样本量持续增大、\(N_{\mathrm{eff}}\) 长期停在个位数、propensity 靠事后模型反推、目标策略频繁选中日志里罕见的动作——出现这些信号时，再调一轮裁剪阈值创造不出缺失的证据。可选的动作是把目标策略往行为策略拉近、只评估被数据充分覆盖的那部分策略、去收一批新数据，或者干脆改报一个保守的下界并说明它保守在哪里。

沿着这条路往下想会碰到一个念头：既然把一个策略的数据掰成另一个策略的估计这么费劲，能不能让被优化的对象直接就是动作分布，用它自己产生的数据算梯度。下一节走的正是这条路，它绕开了值函数里的 \(\max\)，也把新的方差问题一并带来。
